\usection{Planteamiento del Problema}

En la actualidad, la sociedad se encuentra sumergida en una era de información y tecnología que demanda cualidades distintas e innovadoras a los individuos, tanto para su inserción efectiva en el mercado laboral, como para su desenvolvimiento en la cotidianidad. Como indica \citeA{zapataros2015}, en muchos países esto ha derivado en un proceso de formación conocido como alfabetización digital, que puede considerarse una etapa tan fundamental de la educación de los niños y jóvenes como lo son las habilidades de lectura, escritura y matemáticas. Estas capacidades demandadas se asocian tradicionalmente a la capacidad para programar, sin embargo, dicho autor argumenta que existe una forma de pensamiento como trasfondo a esta capacidad, y que tiene un dominio mucho más amplio en cuanto a la resolución de problemas. Se introduce así el concepto de pensamiento computacional (p. 2-3).

\citeA{wing2006} define este término como una aptitud universalmente aplicable, una habilidad fundamental para todos, más allá del área de la computación. Describe el pensamiento computacional como una cualidad que involucra la resolución de problemas, el diseño de sistemas, y el entendimiento del comportamiento humano; aprovechando tareas como la abstracción y la descomposición de problemas grandes y complejos. Profundizando en la aplicación general de esta forma de pensamiento, \citeA{wing2010} señala que la computación ha ganado reconocimiento como el tercer pilar de la ciencia, y que el pensamiento computacional influencia áreas como la medicina, la arqueología, la economía, las finanzas, entre otras; demostrando así su utilidad a nivel general, y su relevancia para ser incluida en programas formativos desde edades tempranas.

Actualmente, la generación de niños y jóvenes ha crecido rodeada de dispositivos y herramientas digitales ubicuas, lo que ha llevado a caracterizarla como una generación de nativos digitales. Sin embargo, la familiaridad con la tecnología no implica, por sí misma, el desarrollo del pensamiento computacional. Es decir, un individuo puede interactuar con dispositivos de forma fluida sin haber construido las estructuras de razonamiento asociadas a la abstracción, la descomposición de problemas o el diseño algorítmico. Esta distinción abre una oportunidad formativa relevante, que resuena con la visión de Seymour Papert, quien en 1971 predecía que las computadoras se volverían accesibles para todos, incluyendo a niños, transformando la forma en que juegan y aprenden \cite{resnick2012}, lo que llama al aprovechamiento pedagógico de las tecnologías. Su enfoque construccionista plantea un aprendizaje en el que las estructuras cognitivas del conocimiento se forman a partir de estructuras físicas y virtuales sobre las cuales los niños puedan reflexionar, iterar y colaborar \cite{nosshoyles2019}.

A pesar del creciente acceso a computadoras, tabletas y teléfonos para personas de todas las edades, \citeA{resnick2012} señala que una parte importante del sueño de Papert sigue sin cumplirse. Los niños principalmente utilizan estos dispositivos para chatear, jugar y buscar información en línea. Sin embargo, una verdadera fluidez digital implicaría que también utilicen estos medios para diseñar y crear sus propios juegos y medios de expresión a través de la programación. Es decir, una parte fundamental del pensamiento computacional tiene que ver con la capacidad de entender las herramientas tecnológicas para resolver problemas asociados a la producción de medios digitales. Aun así, \citeA{resnick2009seeds} destaca que, actualmente, los usos de la tecnología en la educación se enfocan en reforzar métodos antiguos de enseñanza y aprendizaje, en lugar de adaptarse a la era digital, teniendo como consecuencia una deficiencia ante la demanda de profesionales formados con pensamiento computacional.

Tal vez el ejemplo más destacable para crear un entorno de programación que fomente el pensamiento computacional en niños, es Scratch. Según \citeA[p.~62-63]{resnicketal2009}, Scratch se desarrolló en base a las ideas de Papert para que los niños puedan aprender conceptos matemáticos y computacionales, además de ejercitar el pensamiento creativo, racional y sistemático, y el trabajo colaborativo. Para el equipo de Scratch, la capacidad para programar aumenta el rango de formas de crear y expresarse con las computadoras; formando en última instancia el pensamiento computacional a través de habilidades para resolver problemas con estrategias que pueden traducirse a dominios ajenos a la programación. Este lenguaje se diseñó con los principios de piso bajo (es sencillo para empezar), techos altos (da oportunidades para crear proyectos más complejos) y muros anchos (soporte para distintos tipos de proyectos). Scratch se basa en la programación visual, utilizando bloques de instrucciones que permiten codificar programas, orientado a la definición de un flujo de control (control flow).

Si bien se ha tenido éxito utilizando la programación en bloques orientada al control de flujo para la formación del pensamiento computacional, la exploración de herramientas basadas en la definición de flujos de datos en edades tempranas ha sido incipiente y enfocada en dominios específicos. Los lenguajes basados en control flow son efectivos en la enseñanza de la ejecución secuencial, formando lo que Papert llamó \emph{pensamiento procedimental}. Sin embargo, él mismo argumentaba que este era solo un primer paso a la formación del pensamiento computacional, donde el objetivo final es la comprensión y formación de \emph{objetos computacionales} \cite[p.~894]{lodimartini2021}. Para esto, es necesaria la introducción de conceptos más complejos que están cada vez más presentes en la computación moderna, y que pueden aparecer de forma más orgánica en entornos de dataflow.

Estudios recientes han comenzado a evidenciar el potencial de la programación basada en flujos para superar estas barreras. Un ejemplo destacado es M-flow, una plataforma diseñada para enseñar conceptos computacionales a niños a través de la creación musical \cite{mflow2023}. Las investigaciones sobre esta herramienta concluyen que el paradigma de flujos de datos resulta altamente intuitivo, mitigando la frustración asociada a la sintaxis y aumentando significativamente la autoeficacia y el compromiso de los estudiantes hacia la programación. Sin embargo, a pesar de su éxito pedagógico, M-flow presenta una limitación importante, al estar diseñada estrictamente para un dominio específico. Esta restricción a la manipulación de eventos sonoros, hace que carezca de la completitud de Turing y de la versatilidad de un lenguaje de propósito general. A diferencia de entornos como Scratch, que ofrecen ``muros anchos'' \cite{resnicketal2009} para permitir la creación de simulaciones, videojuegos o la resolución de problemas lógicos y matemáticos abstractos, las alternativas de dataflow actuales orientadas a niños no poseen esta capacidad expresiva universal.

Es importante destacar que las herramientas de dataflow pueden ser de gran utilidad en contextos de tecnologías modernas. Un ejemplo de esto es n8n, una herramienta de automatización de flujos de trabajo orientada a tareas basadas en agentes de inteligencia artificial \cite{n8n2025}. Esta herramienta representa un logro importante en el aprovechamiento de IA para personas con pocos conocimientos de programación. De esta manera, se puede identificar una oportunidad para tener un lenguaje con un uso más general, buscando sentar las bases del pensamiento computacional. De manera similar a la que n8n simplifica la interacción con la IA, el lenguaje a desarrollar podría simplificar la comprensión de un conceptos tecnológicos relevantes como la concurrencia y los sistemas distribuidos, partiendo desde las bases del pensamiento computacional, y brindando una enseñanza alternativa para la resolución de problemas, las relaciones causa-efecto y las dependencias de datos para distintos procesos, en lugar de memorizar el orden de ejecución de las operaciones.

Para dar respuesta a la problemática planteada, se propone el desarrollo de un lenguaje de programación visual basado en dataflow y su entorno, diseñado específicamente para ser una herramienta educativa que haga accesible, intuitiva y atractiva la formación del pensamiento computacional en niños y jóvenes de 8 a 16 años.